\documentclass[a4paper,12pt]{article}
\usepackage[utf8]{inputenc}
\usepackage[T1]{fontenc}
\usepackage[ngerman]{babel}
\usepackage{amsmath,amssymb}
\usepackage{booktabs}
\usepackage{enumitem}
\usepackage[margin=2.5cm]{geometry}

\begin{document}

\begin{center}
{\LARGE\bfseries \"Ubungsblatt 3 -- L\"osungen}\\[0.5em]
{\large VC Logik}
\end{center}

\bigskip

%% ============================================================
%% AUFGABE 1
%% ============================================================
\subsection*{Aufgabe 1 \textnormal{(2 Punkte)} -- CNF und DNF f\"ur
$\neg X \leftrightarrow Z \to \neg Y \wedge Z$}

Mit den Pr\"azedenzen $\neg > \wedge > \vee > \to > \leftrightarrow$
klammert sich die Formel als
\[
F \;\equiv\; \neg X \leftrightarrow \bigl(Z \to (\neg Y \wedge Z)\bigr).
\]

Wahrheitstabelle ($2^3 = 8$ Zeilen):

\begin{center}
\renewcommand{\arraystretch}{1.15}
\begin{tabular}{ccc|c|c|c|c}
\toprule
$X$ & $Y$ & $Z$ & $\neg X$ & $\neg Y \wedge Z$
      & $Z \to (\neg Y \wedge Z)$ & $F$ \\
\midrule
0 & 0 & 0 & 1 & 0 & 1 & \textbf{1} \\
0 & 0 & 1 & 1 & 1 & 1 & \textbf{1} \\
0 & 1 & 0 & 1 & 0 & 1 & \textbf{1} \\
0 & 1 & 1 & 1 & 0 & 0 & 0 \\
1 & 0 & 0 & 0 & 0 & 1 & 0 \\
1 & 0 & 1 & 0 & 1 & 1 & 0 \\
1 & 1 & 0 & 0 & 0 & 1 & 0 \\
1 & 1 & 1 & 0 & 0 & 0 & \textbf{1} \\
\bottomrule
\end{tabular}
\end{center}

\medskip
\textbf{Wahre Zeilen:}
$(X,Y,Z) \in \{(0,0,0),\,(0,0,1),\,(0,1,0),\,(1,1,1)\}$.

\textbf{Falsche Zeilen:}
$(X,Y,Z) \in \{(0,1,1),\,(1,0,0),\,(1,0,1),\,(1,1,0)\}$.

\bigskip
\textbf{DNF} (Minterme der wahren Zeilen):
\begin{align*}
\text{DNF}(F) \;\equiv\; &(\neg X \wedge \neg Y \wedge \neg Z) \vee {} \\
& (\neg X \wedge \neg Y \wedge Z) \vee {} \\
& (\neg X \wedge Y \wedge \neg Z) \vee {} \\
& (X \wedge Y \wedge Z).
\end{align*}

\textbf{CNF} (Maxterme der falschen Zeilen, Literale umgekehrt):
\begin{align*}
\text{CNF}(F) \;\equiv\; &(X \vee \neg Y \vee \neg Z) \wedge {} \\
& (\neg X \vee Y \vee Z) \wedge {} \\
& (\neg X \vee Y \vee \neg Z) \wedge {} \\
& (\neg X \vee \neg Y \vee Z).
\end{align*}

\bigskip
\noindent\fbox{\parbox{\dimexpr\textwidth-2\fboxsep-2\fboxrule}{%
\textbf{Ergebnis Aufgabe~1.}\\[2pt]
\textbf{DNF:} $(\neg X \wedge \neg Y \wedge \neg Z)
\vee (\neg X \wedge \neg Y \wedge Z)
\vee (\neg X \wedge Y \wedge \neg Z)
\vee (X \wedge Y \wedge Z)$.\\[2pt]
\textbf{CNF:} $(X \vee \neg Y \vee \neg Z)
\wedge (\neg X \vee Y \vee Z)
\wedge (\neg X \vee Y \vee \neg Z)
\wedge (\neg X \vee \neg Y \vee Z)$.
}}

\newpage

%% ============================================================
%% AUFGABE 2
%% ============================================================
\subsection*{Aufgabe 2 \textnormal{(2 Punkte)} -- Jede Formel ist
\"aquivalent zu einer mit nur $\{\vee,\neg\}$}

\textbf{Behauptung.} Zu jeder aussagenlogischen Formel~$F$ existiert
eine Formel~$F'$, die nur die Junktoren $\vee$ und $\neg$ verwendet
und $F \equiv F'$ erf\"ullt.

\medskip
Beweis per struktureller Induktion \"uber den Formelaufbau.

\medskip
\textbf{Induktionsanfang.} Sei $F = p$ eine Aussagenvariable.
Dann enth\"alt $F$ keine Junktoren, also ist $F' := p$ bereits
eine Formel \"uber $\{\vee, \neg\}$ mit $F' \equiv F$.

\medskip
\textbf{Induktionsschritt.} Sei $F$ zusammengesetzt.
IV: F\"ur alle echten Teilformeln~$G$ von~$F$ existiert ein
\"aquivalentes~$G'$ \"uber $\{\vee, \neg\}$.
Fallunterscheidung nach Hauptoperator:

\medskip
\textbf{Fall 1: $F = \neg G$.}\\
Nach IV existiert $G'$ \"uber $\{\vee,\neg\}$ mit $G \equiv G'$.
Setze $F' := \neg G'$. \checkmark

\medskip
\textbf{Fall 2: $F = G \vee H$.}\\
Nach IV existieren $G', H'$ \"uber $\{\vee, \neg\}$.
Setze $F' := G' \vee H'$. \checkmark

\medskip
\textbf{Fall 3: $F = G \wedge H$.}\\
De~Morgan: $G \wedge H \equiv \neg(\neg G \vee \neg H)$.

$G \wedge H$ ist wahr gdw.\ beide wahr sind, gdw.\ keines falsch ist,
gdw.\ $\neg G \vee \neg H$ falsch ist, gdw.\ $\neg(\neg G \vee \neg H)$
wahr ist. Mit IV:
$F' := \neg(\neg G' \vee \neg H')$. \checkmark

\medskip
\textbf{Fall 4: $F = G \to H$.}\\
$G \to H \equiv \neg G \vee H$
(gleiche Wahrheitstabelle: beide nur falsch wenn $G=1, H=0$).
Mit IV: $F' := \neg G' \vee H'$. \checkmark

\medskip
\textbf{Fall 5: $F = G \leftrightarrow H$.}\\
$G \leftrightarrow H \equiv (G \to H) \wedge (H \to G)$.
Mit Fall~4 (Implikation) und Fall~3 (De~Morgan):
\begin{align*}
G \leftrightarrow H &\;\equiv\; (\neg G \vee H) \wedge (\neg H \vee G) \\
&\;\equiv\; \neg\bigl(\neg(\neg G \vee H) \vee \neg(\neg H \vee G)\bigr).
\end{align*}
Mit IV: $F' := \neg\bigl(\neg(\neg G' \vee H') \vee \neg(\neg H' \vee G')\bigr)$. \checkmark

\medskip
\textbf{Schluss.} Damit ist in jedem Fall eine \"aquivalente Formel
\"uber $\{\vee, \neg\}$ konstruiert. Da jede Formel einen endlichen
Syntaxbaum hat, liefert die Induktion die Behauptung. \hfill$\square$

\bigskip
\noindent\fbox{\parbox{\dimexpr\textwidth-2\fboxsep-2\fboxrule}{%
\textbf{Merke.} $\{\vee, \neg\}$ ist eine \emph{funktional
vollst\"andige} Menge von Junktoren -- man kann jede boolesche
Funktion und damit jede Formel allein mit Disjunktion und Negation
ausdr\"ucken. Die im Beweis verwendeten Ersetzungen sind:
\begin{align*}
G \wedge H &\;\equiv\; \neg(\neg G \vee \neg H) \\
G \to H &\;\equiv\; \neg G \vee H \\
G \leftrightarrow H &\;\equiv\;
   \neg\bigl(\neg(\neg G \vee H) \vee \neg(\neg H \vee G)\bigr)
\end{align*}
}}

\newpage

%% ============================================================
%% AUFGABE 3
%% ============================================================
\subsection*{Aufgabe 3 \textnormal{(2 Punkte)} -- CNF und DNF aus
Wahrheitstabelle \"uber $\sigma = \{A,B,C,D\}$}

Wahrheitstabelle mit $2^4 = 16$ Zeilen:

\begin{center}
\renewcommand{\arraystretch}{1.1}
\begin{tabular}{r|cccc|c|l}
\toprule
\# & $A$ & $B$ & $C$ & $D$ & $f$ & verwendet f\"ur \\
\midrule
 1 & 0 & 0 & 0 & 0 & 0 & CNF-Maxterm \\
 2 & 0 & 0 & 0 & 1 & \textbf{1} & DNF-Minterm \\
 3 & 0 & 0 & 1 & 0 & \textbf{1} & DNF-Minterm \\
 4 & 0 & 0 & 1 & 1 & 0 & CNF-Maxterm \\
 5 & 0 & 1 & 0 & 0 & \textbf{1} & DNF-Minterm \\
 6 & 0 & 1 & 0 & 1 & \textbf{1} & DNF-Minterm \\
 7 & 0 & 1 & 1 & 0 & 0 & CNF-Maxterm \\
 8 & 0 & 1 & 1 & 1 & \textbf{1} & DNF-Minterm \\
 9 & 1 & 0 & 0 & 0 & 0 & CNF-Maxterm \\
10 & 1 & 0 & 0 & 1 & 0 & CNF-Maxterm \\
11 & 1 & 0 & 1 & 0 & \textbf{1} & DNF-Minterm \\
12 & 1 & 0 & 1 & 1 & 0 & CNF-Maxterm \\
13 & 1 & 1 & 0 & 0 & \textbf{1} & DNF-Minterm \\
14 & 1 & 1 & 0 & 1 & 0 & CNF-Maxterm \\
15 & 1 & 1 & 1 & 0 & \textbf{1} & DNF-Minterm \\
16 & 1 & 1 & 1 & 1 & \textbf{1} & DNF-Minterm \\
\bottomrule
\end{tabular}
\end{center}

\medskip
9~wahre Zeilen $\to$ 9~Minterme (DNF),
7~falsche Zeilen $\to$ 7~Maxterme (CNF).

\bigskip
\textbf{DNF} (Minterme der wahren Zeilen):

\begin{align*}
\text{DNF}(f) \;\equiv\; &(\neg A \wedge \neg B \wedge \neg C \wedge D) \vee {} \\
& (\neg A \wedge \neg B \wedge C \wedge \neg D) \vee {} \\
& (\neg A \wedge B \wedge \neg C \wedge \neg D) \vee {} \\
& (\neg A \wedge B \wedge \neg C \wedge D) \vee {} \\
& (\neg A \wedge B \wedge C \wedge D) \vee {} \\
& (A \wedge \neg B \wedge C \wedge \neg D) \vee {} \\
& (A \wedge B \wedge \neg C \wedge \neg D) \vee {} \\
& (A \wedge B \wedge C \wedge \neg D) \vee {} \\
& (A \wedge B \wedge C \wedge D).
\end{align*}

\textbf{CNF} (ein Maxterm pro falscher Zeile, insgesamt 7~St\"uck):

\begin{align*}
\text{CNF}(f) \;\equiv\; &(A \vee B \vee C \vee D) \wedge {} \\
& (A \vee B \vee \neg C \vee \neg D) \wedge {} \\
& (A \vee \neg B \vee \neg C \vee D) \wedge {} \\
& (\neg A \vee B \vee C \vee D) \wedge {} \\
& (\neg A \vee B \vee C \vee \neg D) \wedge {} \\
& (\neg A \vee B \vee \neg C \vee \neg D) \wedge {} \\
& (\neg A \vee \neg B \vee C \vee \neg D).
\end{align*}

\medskip
\medskip
\textbf{Stichprobe:} Zeile~9, $(A,B,C,D)=(1,0,0,0)$, $f=0$.
Maxterm $(\neg A \vee B \vee C \vee D) = (0 \vee 0 \vee 0 \vee 0) = 0$
$\Rightarrow$ CNF wird in dieser Zeile~$0$. \checkmark

\end{document}
